\section{Gysin maps and push--pull operators}\label{sec:pushpull}

The cohomology ring presentation of Subsection~\ref{ss:equiv_coho_xpsi} provides
an explicit polynomial description of the Gysin maps for lattice inclusions
and the push--pull operators for forgetful morphisms.
Throughout this section, we work integrally over the coefficient ring $\Lambda_K$,
with $K=T\times S$ or $K=G\times S$, and retain the notation
$R_K=\Lambda_K[x_1,\dots,x_n]$, $f_K$, and $I_{\rfunc,K}$, with $x_i=-c_1^K(L_i/L_{i-1})$.

For $\rfunc'\leq\rfunc$, the tautological line bundles on $\SX_{\rfunc}$ restrict along the inclusion $j:\SX_{\rfunc'}\hookrightarrow\SX_{\rfunc}$ to those on $\SX_{\rfunc'}$. Consequently, $I_{\rfunc,K}\subset I_{\rfunc',K}$ and the pullback $j^*$ is identified with the natural quotient map $j^*:R_K/I_{\rfunc,K}\twoheadrightarrow R_K/I_{\rfunc',K}$.

\subsection{Gysin maps for divisor embeddings}

Let $\rfunc$ and $\rfunc'$ be reverse Hessenberg functions such that $\rfunc'$ is obtained
from $\rfunc$ by decrementing a single value:
\[
  \rfunc'(m)=\rfunc(m)-1 \quad\text{for some } m\in\{2,\dots,n\},
  \qquad \rfunc'(k)=\rfunc(k) \quad(k\neq m).
\]
Set $a=\rfunc(m)$ and let $i:\SX_{\rfunc'}\hookrightarrow\SX_{\rfunc}$
be the canonical inclusion.

\begin{proposition}\label{prop:inclusion_gysin}
The closed embedding $i:\SX_{\rfunc'}\hookrightarrow\SX_{\rfunc}$ realizes
$\SX_{\rfunc'}$ as a smooth Cartier divisor in $\SX_{\rfunc}$, with normal bundle
isomorphic to the restriction of the equivariant line bundle
\[
  \mathcal E=\Hom(tL_m/tL_{m-1},L_a/L_{a-1}).
\]
Under the presentations of Theorem~\ref{thm:equiv_coho_xdelta} and Corollary~\ref{cor:equiv_coho_xdelta_G}, the Gysin pushforward map
$i_*:H_K^q(\SX_{\rfunc'};\Z)\to H_K^{q+2}(\SX_{\rfunc};\Z)$ is given by
\[
  i_*[f]=[(x_m-x_a-\epsilon)f] \qquad (f\in R_K).
\]
\end{proposition}

\begin{proof}
Weak monotonicity of $\rfunc'$ gives $\rfunc(m-1)=\rfunc'(m-1)\leq\rfunc'(m)=a-1$,
so $tL_{m-1}\subset L_{a-1}$ holds everywhere on $\SX_{\rfunc}$.
Together with $tL_m\subset L_a$, the lattice inclusion induces an equivariant homomorphism
$tL_m/tL_{m-1}\to L_a/L_{a-1}$, defining a section $s\in\Gamma(\SX_{\rfunc},\mathcal E)$
whose zero locus is precisely $\SX_{\rfunc'}$.
At stage $m$ of the projective bundle tower, $tL_{m-1}\subset L_{a-1}$
makes $t^{-1}L_{a-1}/L_{m-1}$ a corank-one subbundle of
$t^{-1}L_a/L_{m-1}$, and $s=0$ defines the smooth relative hyperplane
\[
 \PP(t^{-1}L_{a-1}/L_{m-1})\subset\PP(t^{-1}L_a/L_{m-1}).
\]
Let $p_m:\SX_{\rfunc}\to\SX_{\rfunc}^{[m]}$ forget the lattices after $L_m$.
Since only the $m$-th lattice condition changes, $\SX_{\rfunc'}$ is the
scheme-theoretic inverse image of this hyperplane under $p_m$.
The map $p_m$ is a composition of projective bundle projections and is therefore smooth.
Thus $s$ is transverse to the zero section and the normal bundle is $i^*\mathcal E$.

Since multiplication by $t$ twists by $\C_1$, the first Chern classes of the two quotient lines are $c_1^K(tL_m/tL_{m-1})=-x_m+\epsilon$ and $c_1^K(L_a/L_{a-1})=-x_a$, giving
$c_1^K(\mathcal E)=x_m-x_a-\epsilon$.
Because $i_*(1)=c_1^K(\mathcal E)$, the projection formula yields $i_*(i^*\alpha)=\alpha\cup c_1^K(\mathcal E)=(x_m-x_a-\epsilon)\alpha$
for all $\alpha\in H_K^*(\SX_{\rfunc};\Z)$.
Since $i^*$ is surjective, every class on $\SX_{\rfunc'}$ is of the form $i^*[f]=[f]$,
yielding $i_*[f]=[(x_m-x_a-\epsilon)f]$.
\end{proof}

\subsection{Forgetting a lattice}

We next consider the forgetful morphism that drops the intermediate lattice $L_m$ from a flag.
Fix an index $1\leq m<n$, and assume that
\begin{equation}\label{eq:pushpull_hypotheses}
  \rfunc(m)=\rfunc(m+1),\qquad
  \rfunc(k)\neq m\quad(k\geq m+2).
\end{equation}
Let $K_m=\{0,\dots,n\}\setminus\{m\}$.
By \eqref{eq:pushpull_hypotheses}, $\rfunc(k)\neq m$ for every $k$,
so all targets $L_{\rfunc(k)}$ belong to the retained flag.
We may therefore define the partial flag variety
\[
  \SX_{\rfunc,m}^P
  =\left\{(L_k)_{k\in K_m}\ \middle|\;
  \begin{gathered}
    L_k\in\Gr,\quad L_0=t\CO^n,\quad L_k\subset L_\ell\ (k<\ell),\\
    \dim_\C(L_k/L_0)=k,\quad
    tL_k\subset L_{\rfunc(k)}\ (k\in K_m\setminus\{0\})
  \end{gathered}\right\},
\]
and denote by $\pi_m:\SX_{\rfunc}\to\SX_{\rfunc,m}^P$ the morphism forgetting $L_m$.

\begin{proposition}\label{prop:pushpull_projective_bundle}
Under \eqref{eq:pushpull_hypotheses}, the morphism $\pi_m$ realizes $\SX_{\rfunc}$ as the
equivariant projective bundle $\PP(\mathcal V_m)$ over $\SX_{\rfunc,m}^P$, where
$\mathcal V_m=L_{m+1}/L_{m-1}$ is a rank-two equivariant vector bundle.
In particular, the base variety $\SX_{\rfunc,m}^P$ is smooth and projective.
\end{proposition}

\begin{proof}
Set $a=\rfunc(m)=\rfunc(m+1)\leq m-1$. For any partial flag in $\SX_{\rfunc,m}^P$,
we have $tL_{m+1}\subset L_a\subset L_{m-1}$, so every intermediate lattice
$L_{m-1}\subset L_m\subset L_{m+1}$ automatically satisfies $tL_m\subset tL_{m+1}\subset L_a$.
Since no other condition refers to $L_m$ by \eqref{eq:pushpull_hypotheses},
a lattice $L_m$ completes the partial flag if and only if $L_{m-1}\subset L_m\subset L_{m+1}$ with $\dim_\C(L_m/L_{m-1})=1$.
Such lattices correspond bijectively to lines in $L_{m+1}/L_{m-1}$,
identifying $\SX_{\rfunc}\cong\PP(\mathcal V_m)$ equivariantly.
The base is smooth because the projective bundle is locally a product with $\PP^1$
and $\SX_{\rfunc}$ is smooth.
For projectivity, the $n$-dimensional $\CO$-module $L_n/L_0$ is annihilated by $t^n$,
so every retained subspace $L_k/L_0$ lies in the fixed vector space $t^{-n}L_0/L_0$.
The incidence and lattice conditions are closed conditions in the corresponding
partial flag variety, so $\SX_{\rfunc,m}^P$ is projective.
\end{proof}

\subsection{The geometric divided difference operator}

We now relate the push--pull operator $\pi_m^*\pi_{m*}$ to the divided difference
operator $\partial_m$ on $R_K$.

\begin{theorem}\label{thm:pushpull_divided_difference}
Assume \eqref{eq:pushpull_hypotheses}.
Let $\Phi_K:R_K\twoheadrightarrow H_K^*(\SX_{\rfunc};\Z)$ be the presentation homomorphism.
Then for every $f\in R_K$,
\[
  \pi_m^*\pi_{m*}\Phi_K(f)=\Phi_K(\partial_m f).
\]
Consequently, $I_{\rfunc,K}$ is stable under $\partial_m$ and $\tau_m$.
\end{theorem}

\begin{proof}
On $\SX_{\rfunc}$, the pulled-back bundle $\pi_m^*\mathcal V_m=L_{m+1}/L_{m-1}$ fits into
the tautological exact sequence $0\to L_m/L_{m-1}\to\pi_m^*\mathcal V_m\to L_{m+1}/L_m\to 0$.
Taking equivariant Chern classes yields $c_u^K(\pi_m^*\mathcal V_m)=(1-x_m u)(1-x_{m+1}u)$, so
\[
  \Phi_K(x_m+x_{m+1})=-\pi_m^*c_1^K(\mathcal V_m),\qquad
  \Phi_K(x_m x_{m+1})=\pi_m^*c_2^K(\mathcal V_m).
\]
For every index $k\notin\{m,m+1\}$, the line bundle $L_k/L_{k-1}$ is pulled back from $\SX_{\rfunc,m}^P$,
so $\Phi_K(x_k)\in\pi_m^*H_K^2(\SX_{\rfunc,m}^P;\Z)$.

Let $R_K^{s_m}\subset R_K$ be the subring of polynomials symmetric in $x_m,x_{m+1}$.
Every $f\in R_K$ can be written uniquely as $f=a+x_m b$ with $a,b\in R_K^{s_m}$,
whence $\partial_m f=b$ and both $\Phi_K(a)$ and $\Phi_K(b)$ lie in $\pi_m^*H_K^*(\SX_{\rfunc,m}^P;\Z)$.
Since $\pi_{m*}1=0$ and $\pi_{m*}x_m=1$, the projection formula gives
$\pi_m^*\pi_{m*}\Phi_K(f)=\Phi_K(b)=\Phi_K(\partial_m f)$.

Since $I_{\rfunc,K}=\ker\Phi_K$, if $f\in I_{\rfunc,K}$ then
$\Phi_K(\partial_m f)=\pi_m^*\pi_{m*}(0)=0$, so $\partial_m f\in I_{\rfunc,K}$.
From $\tau_m=\id-(x_m-x_{m+1})\partial_m$, it follows that $I_{\rfunc,K}$ is also stable under $\tau_m$.
\end{proof}

\subsection{Gysin maps for arbitrary inclusions}

We now generalize the Gysin formula of Proposition~\ref{prop:inclusion_gysin} to arbitrary inclusions.
Let $\rfunc'\leq\rfunc$ be reverse Hessenberg functions, and let
$j:\SX_{\rfunc'}\hookrightarrow\SX_{\rfunc}$ be the canonical closed embedding of codimension
$c=\sum_{k=1}^n(\rfunc(k)-\rfunc'(k))$.
Define the polynomial
\[
  E_{\rfunc',\rfunc}
  =\prod_{k=1}^n\prod_{a=\rfunc'(k)+1}^{\rfunc(k)}
       (x_k-x_a-\epsilon)\ \in R_K,
\]
which is homogeneous of cohomological degree $2c$.

\begin{proposition}\label{prop:general_gysin}
The Gysin pushforward $j_*:H_K^q(\SX_{\rfunc'};\Z)\to H_K^{q+2c}(\SX_{\rfunc};\Z)$ is given by
\[
  j_*[f]=[E_{\rfunc',\rfunc}f] \qquad (f\in R_K).
\]
Moreover, $j_*$ is injective.
\end{proposition}

\begin{proof}
We transition from $\rfunc$ to $\rfunc'$ by decreasing the value at index $k$ from $\rfunc(k)$ down to $\rfunc'(k)$ one unit at a time, proceeding in increasing order $k=1,\dots,n$.
At step $k$, the values at indices $h<k$ have already been reduced to $\rfunc'(h)$,
while those at $h>k$ remain $\rfunc(h)$.
Any intermediate value $v$ with $\rfunc'(k)\leq v\leq\rfunc(k)$ satisfies $0\leq v<k$ and
\[
 \begin{aligned}
  \rfunc'(k-1)&\leq\rfunc'(k)\leq v &&\text{if }k>1,\\
  v&\leq\rfunc(k)\leq\rfunc(k+1) &&\text{if }k<n.
 \end{aligned}
\]
Thus each intermediate function remains a reverse Hessenberg function.
This factors $j$ into $c$ inclusions of smooth Cartier divisors.
By Proposition~\ref{prop:inclusion_gysin}, each decrement from $a$ to $a-1$
at index $k$ contributes the factor $x_k-x_a-\epsilon$.
Functoriality of Gysin maps therefore gives $j_*[f]=[E_{\rfunc',\rfunc}f]$.

If $j_*\alpha=0$, the projection formula gives
\[
 \int_{\SX_{\rfunc'}}\alpha\cup j^*\beta
 =\int_{\SX_{\rfunc}}j_*\alpha\cup\beta=0
 \qquad\bigl(\beta\in H_K^*(\SX_{\rfunc};\Z)\bigr).
\]
Since $j^*$ is surjective and the integration pairing on $\SX_{\rfunc'}$
is perfect, we conclude that $\alpha=0$.
\end{proof}
